\begin{lemma}
\label{lemma-characterize-K-injective}
Let $\mathcal{A}$ be a Grothendieck abelian category.
Let $\kappa$ be a cardinal as in
Lemma \ref{lemma-acyclic-quotient-complexes-bounded-size}.
Suppose that $I^\bullet$ is a complex such that
\begin{enumerate}
\item each $I^j$ is injective, and
\item for every bounded above acyclic complex $M^\bullet$
such that $|M^n| \leq \kappa$
we have $\Hom_{K(\mathcal{A})}(M^\bullet, I^\bullet) = 0$.
\end{enumerate}
Then $I^\bullet$ is an $K$-injective complex.
\end{lemma}

\begin{proof}
Let $M^\bullet$ be an acyclic complex. We are going to construct by
induction on the ordinal $\alpha$ an acyclic subcomplex
$K_\alpha^\bullet \subset M^\bullet$ as follows.
For $\alpha = 0$ we set $K_0^\bullet = 0$. For $\alpha > 0$
we proceed as follows:
\begin{enumerate}
\item If $\alpha = \beta + 1$ and $K_\beta^\bullet = M^\bullet$
then we choose $K_\alpha^\bullet = K_\beta^\bullet$.
\item If $\alpha = \beta + 1$ and $K_\beta^\bullet \not = M^\bullet$
then $M^\bullet/K_\beta^\bullet$ is a nonzero acyclic complex.
We choose a subcomplex $N_\alpha^\bullet \subset M^\bullet/K_\beta^\bullet$
as in Lemma \ref{lemma-acyclic-quotient-complexes-bounded-size}.
Finally, we let $K_\alpha^\bullet \subset M^\bullet$
be the inverse image of $N_\alpha^\bullet$.
\item If $\alpha$ is a limit ordinal we set
$K_\beta^\bullet = \colim K_\alpha^\bullet$.
\end{enumerate}
It is clear that $M^\bullet = K_\alpha^\bullet$ for a suitably large
ordinal $\alpha$. We will prove that
$$
\Hom_{K(\mathcal{A})}(K_\alpha^\bullet, I^\bullet)
$$
is zero by transfinite induction on $\alpha$. It holds for $\alpha = 0$
since $K_0^\bullet$ is zero. Suppose it holds for $\beta$ and
$\alpha = \beta + 1$. In case (1) of the list above the result is clear.
In case (2) there is a short exact sequence of complexes
$$
0 \to K_\beta^\bullet \to K_\alpha^\bullet \to N_\alpha^\bullet \to 0
$$
Since each component of $I^\bullet$ is injective we see that we obtain
an exact sequence
$$
\Hom_{K(\mathcal{A})}(K_\beta^\bullet, I^\bullet) \leftarrow
\Hom_{K(\mathcal{A})}(K_\alpha^\bullet, I^\bullet) \leftarrow
\Hom_{K(\mathcal{A})}(N_\alpha^\bullet, I^\bullet)
$$
By induction the term on the left is zero and by assumption on $I^\bullet$
the term on the right is zero. Thus the middle group is zero too.
Finally, suppose that $\alpha$ is a limit ordinal. Then we see that
$$
\Hom^\bullet(K_\alpha^\bullet, I^\bullet) =
\lim_{\beta < \alpha} \Hom^\bullet(K_\beta^\bullet, I^\bullet)
$$
with notation as in
More on Algebra, Section \ref{more-algebra-section-hom-complexes}.
These complexes compute morphisms in $K(\mathcal{A})$ by
More on Algebra, Equation
(\ref{more-algebra-equation-cohomology-hom-complex}).
Note that the transition maps in the system are surjective
because $I^j$ is injective for each $j$. Moreover, for a limit
ordinal $\alpha$ we have equality of limit and value
(see displayed formula above). Thus we may apply
Homology, Lemma \ref{homology-lemma-ML-over-ordinals}
to conclude.
\end{proof}
